% paper/sections/theorem1_bound.tex -- P2-02, plan/03-phase2-theory.md.
%
% Standalone section fragment, \input{} by paper/main.tex once P6-01 builds
% the skeleton. Requires the same preamble items as setup.tex (amsmath,
% amssymb, amsthm, mathtools, theorem/definition/remark environments,
% \needshuman{...}), plus \newtheorem{corollary}{Corollary} and
% \newtheorem{lemma}{Lemma}. Depends on paper/sections/setup.tex's notation
% (Sections~\ref{sec:setup}--\ref{sec:setup-policy}).
%
% IMPORTANT for whoever reads or integrates this: the bound stated here is
% NOT the form plan/03-phase2-theory.md's draft text states
% (`exp(-Delta_k^2/(2*(sigma2_extrap_k+v_k)))`, additive in
% sigma^2_extrap), and not the form already implemented as
% `marginal_bound_term`/`pairwise_bound_term` in src/pdt/theory/bound.py
% as first shipped (used in the first P1-07/P1-08 runs; since corrected to the
% gap-reduction form by PR #17). That additive form was checked with
% tests/theory/test_theorem1.py's numerical certificate and found to be a
% genuinely INVALID worst-case bound -- see the remark after the proof, and
% docs/decisions.md, 2026-09-10, for the full account of how this was
% found and why P1-07/P1-08's already-reported numbers are not affected by
% the correction (they describe a real, empirically-well-behaved-on-real-
% data formula; what changes is only the claim that the formula is a
% *proven* worst-case bound in general).

\section{Theorem 1: the extrapolation-aware error bound}
\label{sec:theorem1}

\subsection{Statement}

\begin{theorem}[Extrapolation-aware selection-error bound]
\label{thm:main-bound}
Fix a design $S_{\mathrm{fit}}$ with total compute $C$. Arm $k$'s true trend is
$\mu_k(s) = g(\theta_k, s) + h_k(s)$ with $h_k \in \mathcal{H}$
(Assumption~\ref{asn:misspec}); $\hat\mu_k(s^\star) = g(\hat\theta_k, s^\star)$, with
$\hat\theta_k$ the (weighted) least-squares fit within the family
(Assumption~\ref{asn:family}) on noisy observations (Assumption~\ref{asn:noise}), and the
projection $\theta_k^{\dagger}$ of Definition~\ref{def:projection} taken with $\pi$ the
design's own weighting. Write $\sigma^2_{\mathrm{extrap}, k}
= \sigma^2_{\mathrm{extrap}, k}(S_{\mathrm{fit}})$ (Definition~\ref{def:sigma2-extrap}),
$\Delta_k = \mu_{k^\star}(s^\star) - \mu_k(s^\star)$, and let
$v_k(C) = J(\theta_k^{\dagger}, s^\star)^\top \Sigma_{\theta_k}(C) \, J(\theta_k^{\dagger},
s^\star)$ be the delta-method variance (Section~\ref{sec:theorem1-step2}). Define
\[
  B(\rho_k, \rho_{k^\star}, v_k, v_{k^\star}) \;=\; \sum_{k \ne k^\star}
  \exp\!\left(
    -\, \frac{\Big( \Delta_k - \sqrt{\sigma^2_{\mathrm{extrap}, k}}
      - \sqrt{\sigma^2_{\mathrm{extrap}, k^\star}} - \rho_k - \rho_{k^\star} \Big)_+^2}
    {2 \big( v_k + v_{k^\star} \big)}
  \right),
  \qquad (x)_+ = \max(x, 0).
\]

\textbf{(i) Linear-in-$\theta$ family, independent sub-Gaussian noise (proved).} If
$g(\theta, s)$ is linear in $\theta$ and the observation noise is independent
sub-Gaussian with variance proxy $\sigma^2(s)$ at each fitting observation, the fit is
a fixed linear function of the noise, $\hat\mu_k(s^\star)$ is exactly sub-Gaussian
around $g(\theta_k^{\dagger}, s^\star)$ with variance proxy $v_k(C)$, and
\[
  \mathbb{P}[\hat{k} \ne k^\star] \;\le\; B(0, 0, v_k(C), v_{k^\star}(C)).
\]
\textbf{(ii) Nonlinear $g$ (conditional, not proved here).} For nonlinear $g$ the
delta-method variance $v_k(C)$ is a first-order approximation and the fitted prediction
is neither exactly sub-Gaussian nor centred at $g(\theta_k^{\dagger}, s^\star)$: the
Taylor remainder and the finite-sample bias of the nonlinear least-squares estimator
shift its centre by some $\rho_k(C)$, and its true tails are governed by some proxy
$\bar v_k(C)$ that need not equal $v_k(C)$. \emph{If} one assumes
\begin{quote}
\textbf{(H)} $\hat\mu_k(s^\star) - g(\theta_k^{\dagger}, s^\star)$ is sub-Gaussian with
mean of magnitude at most $\rho_k(C)$ and variance proxy $\bar v_k(C)$, for every arm,
independently across arms,
\end{quote}
\emph{then} $\mathbb{P}[\hat{k} \ne k^\star] \le B(\rho_k, \rho_{k^\star}, \bar v_k,
\bar v_{k^\star})$. Hypothesis (H) is not established by smoothness or a bounded
Jacobian; it is stated as a hypothesis, and the six fitters of
\texttt{src/pdt/scaling/fitters.py} are treated as satisfying it only heuristically
(P1-07 found the delta-method $v_k$ close to the bootstrap variance on real cells, which
is evidence about the variance proxy, not about (H)'s concentration or about $\rho_k$).
Every reported use of Theorem~\ref{thm:main-bound} with a nonlinear fitter --
\texttt{pdt.theory.bound}, P1-07/P1-08 -- is therefore an \emph{empirical diagnostic
built on this conditional statement}, with $\rho_k = 0$ and $\bar v_k = v_k$.

As $C \to \infty$ with the design shape (the relative weighting of $S_{\mathrm{fit}}$
across scales) held fixed, $v_k(C) \to 0$ for every $k$, while
$\sigma^2_{\mathrm{extrap}, k}(S_{\mathrm{fit}})$ does not depend on $C$ at all.
\end{theorem}

\begin{remark}[This is not the additive form]
\label{rem:not-additive}
An earlier draft of this bound (and \texttt{src/pdt/theory/bound.py}'s originally shipped
\texttt{marginal\_bound\_term}, used in the first P1-07/P1-08 runs; PR~\#17 has since corrected
the plug-in to the gap-reduction form) instead adds
$\sigma^2_{\mathrm{extrap}, k}$ directly into the denominator,
$\exp(-\Delta_k^2 / (2(\sigma^2_{\mathrm{extrap}, k} + v_k)))$, using only arm $k$'s own
quantities. \textbf{This additive form is not a valid worst-case bound}: a concrete
counter-example (a fixed bias of $10$ against a gap of $5$ and noise variance $0.01$) %NUMBER-OK: hand-chosen illustrative constants, also in tests/test_bound.py::test_plug_in_and_worst_case_forms_both_reject_the_additive_counterexample -- not a reported result
gives a true error probability of essentially $1$ against a claimed bound of $0.88$, %NUMBER-OK: same hand-chosen counter-example, computed value not a reported result
a genuine violation, confirmed both by hand and by
\texttt{tests/theory/test\_theorem1.py}'s 5000-instance numerical certificate before
either form was trusted. The correct worst-case treatment of a \emph{fixed but
unknown-sign} bias bounded in magnitude by $\sqrt{\sigma^2_{\mathrm{extrap}, k}}$
subtracts that magnitude from the gap (a Chernoff bound evaluated at the adversarial
bias), rather than adding its square to the variance -- which is also why $k^\star$'s
own bias and variance must appear symmetrically: the compared quantity is
$\hat\mu_k(s^\star) - \hat\mu_{k^\star}(s^\star)$, and both sides are noisy estimates.
See \texttt{docs/decisions.md} (2026-09-10) for the full derivation, the numerical
evidence, and why P1-07's ``0 violations on 396 real cells'' finding does not
contradict this: every real DataDecide cell's bound value was $\ge 1$ (vacuous) either
way, so the two forms were never actually distinguished by that check.
\end{remark}

\subsection{Proof}

\paragraph{Step 1 -- concentration of the least-squares estimate.}
\label{sec:theorem1-step1}
For a \emph{linear-in-$\theta$} family (Assumption~\ref{asn:family} with $g(\theta, s)$
linear in $\theta$), $\hat\theta_k$ is an exact linear function of the noise
$\{\varepsilon_{k,i}(s)\}$, hence exactly Gaussian (if $\varepsilon$ is Gaussian) or
exactly sub-Gaussian with the stated parameter (for general sub-Gaussian $\varepsilon$,
via sub-Gaussian linearity) -- no approximation is involved, and
Step~\ref{sec:theorem1-step2}'s delta method is exact rather than a linearization. This
is the whole of Theorem~\ref{thm:main-bound}(i)'s proof of concentration.

For a \emph{nonlinear} $g$ this step is \textbf{not proved}. $\hat\theta_k$ is an
M-estimator minimizing a (noisy, misspecified) sum of squares; consistency and
asymptotic normality as $C \to \infty$ are standard under smoothness and a non-degenerate
Hessian of the population objective near $\theta_k^{\dagger}$ (van der Vaart, 1998,
Ch.~5), but an asymptotic normal limit gives neither a non-asymptotic sub-Gaussian tail
nor a finite-$C$ bound on the centring error $\rho_k(C)$ (Taylor remainder plus
finite-sample estimator bias), and smoothness with a bounded Jacobian does not supply
them. Theorem~\ref{thm:main-bound}(ii) therefore takes exactly that concentration
statement as hypothesis (H) rather than deriving it.
\needshuman{Making (H) a theorem for a concrete nonlinear family means deriving explicit
non-asymptotic constants (local Hessian condition number, Lipschitz constants,
$\rho_k(C)$, $\bar v_k(C)$) valid at finite $C$. That is genuine technical work, not a
citation lookup, and is deliberately not attempted here.}

\paragraph{Step 2 -- delta-method propagation to $s^\star$.}
\label{sec:theorem1-step2}
Mechanical, given Step 1: $\hat\mu_k(s^\star) - g(\theta_k^\dagger, s^\star)
= g(\hat\theta_k, s^\star) - g(\theta_k^\dagger, s^\star) \approx
J(\theta_k^\dagger, s^\star)^\top (\hat\theta_k - \theta_k^\dagger)$ (exact for linear
$g$; a first-order Taylor remainder, controlled by Assumption~\ref{asn:family}'s bounded
Jacobian, for nonlinear $g$), giving $v_k(C) = J^\top \Sigma_{\theta_k}(C) \, J$ as
claimed. $\Sigma_{\theta_k}(C) = (X^\top X)^{-1}(X^\top \mathrm{diag}(\sigma^2) X)(X^\top
X)^{-1}$ for design matrix $X$ (one row $\partial g/\partial\theta$ per fitting
observation) is exactly \texttt{sandwich\_covariance()} /
\texttt{analytic\_v\_k()} in \texttt{src/pdt/theory/bound.py}, already implemented and
tested (P1-07). $v_k(C) = O(1/n_{\mathrm{replicates}}) = O(1/C)$ for a fixed design
shape scaled up by more replicates: mechanical, since $\Sigma_{\theta_k}(C)$'s
`meat'-and-`bread' sandwich structure scales the estimating-equation variance by
$1/n_{\mathrm{replicates}}$ while $X$ (hence $J$) is unchanged --
\texttt{tests/theory/test\_theorem1.py::test\_theorem1\_v\_k\_shrinks\_with\_replicates}
checks this directly (strictly decreasing, and shrinks by $\ge 10^4\times$ as
replicates grow $10^4\times$, exactly as the closed form predicts).
$\sigma^2_{\mathrm{extrap}, k}(S_{\mathrm{fit}})$'s independence from $C$ is immediate
from Definition~\ref{def:projection}: $\theta_k^\dagger$ is defined by an expectation
over the design measure $\pi$ alone, with no dependence on how many replicates are
drawn at each point in its support.

\paragraph{Step 3 -- bias-plus-deviation decomposition.}
\label{sec:theorem1-step3}
Write $Z_k = \hat\mu_k(s^\star) - \mu_k(s^\star)$, so $Z_k$ has (population) mean
$\mathrm{bias}_k := g(\theta_k^\dagger, s^\star) - \mu_k(s^\star)$, with
$\mathrm{bias}_k^2 = \sigma^2_{\mathrm{extrap}, k}$, and is (by Steps 1--2 in case~(i); by hypothesis (H) in case~(ii), where
$\mathrm{bias}_k$ is shifted by at most $\rho_k$ and $v_k$ replaced by $\bar v_k$)
sub-Gaussian around $\mathrm{bias}_k$ with parameter $v_k(C)$. The event
$\{\hat{k} \ne k^\star\}$ is contained in $\bigcup_{k \ne k^\star}
\{\hat\mu_k(s^\star) \ge \hat\mu_{k^\star}(s^\star)\}
= \bigcup_{k \ne k^\star} \{Z_k - Z_{k^\star} \ge \Delta_k\}$. For fixed $k$, since
$Z_k - Z_{k^\star}$ is sub-Gaussian with mean $\mathrm{bias}_k - \mathrm{bias}_{k^\star}$
and parameter $v_k(C) + v_{k^\star}(C)$ (sub-Gaussian additivity under the independence
of arm $k$'s and $k^\star$'s fits -- a modelling choice; Corollary~\ref{cor:pairwise}
below removes it), a one-sided Chernoff bound gives, for any fixed
$\mathrm{bias}_k, \mathrm{bias}_{k^\star}$ with $|\mathrm{bias}_k| \le
\sqrt{\sigma^2_{\mathrm{extrap},k}}$, $|\mathrm{bias}_{k^\star}| \le
\sqrt{\sigma^2_{\mathrm{extrap},k^\star}}$:
\[
  \mathbb{P}[Z_k - Z_{k^\star} \ge \Delta_k]
  \le \exp\!\left(
    - \frac{\big(\Delta_k - (\mathrm{bias}_k - \mathrm{bias}_{k^\star})\big)_+^2}
    {2(v_k(C) + v_{k^\star}(C))} \right).
\]
Taking the worst case over the unknown signs of $\mathrm{bias}_k, \mathrm{bias}_{k^\star}$
(an adversary sets $\mathrm{bias}_k = +\sqrt{\sigma^2_{\mathrm{extrap},k}}$,
$\mathrm{bias}_{k^\star} = -\sqrt{\sigma^2_{\mathrm{extrap},k^\star}}$, minimising the
numerator) gives exactly the per-term bound in the theorem statement. Mechanical, given
Step 1's concentration and the standard one-sided Gaussian/sub-Gaussian tail inequality
$\mathbb{P}[X - \mathbb{E}X \ge t] \le \exp(-t^2/(2s^2))$ for $t \ge 0$.

\paragraph{Step 4 -- union bound.}
\label{sec:theorem1-step4}
$\mathbb{P}[\hat{k}\ne k^\star] = \mathbb{P}\big[\bigcup_{k\ne k^\star}\{\cdot\}\big]
\le \sum_{k \ne k^\star} \mathbb{P}[\cdot]$ by subadditivity -- mechanical, no
independence or further structure required. $\blacksquare$

\subsection{Corollary 1 -- consistency as \texorpdfstring{$C \to \infty$}{C to infinity}:
a strict pairwise-ordering condition is sufficient}
\label{sec:corollary1}

\begin{definition}[Pairwise difference statistic and its population bias]
\label{def:pairwise-d}
$D_k = \hat\mu_{k^\star}(s^\star) - \hat\mu_k(s^\star)$. Its population (noise-free,
$C \to \infty$) limit is $D_k^\dagger = g(\theta_{k^\star}^\dagger, s^\star) -
g(\theta_k^\dagger, s^\star)$, and $\mathrm{bias}(D_k) = D_k^\dagger - \Delta_k$.
\end{definition}

\begin{corollary}[Asymptotic consistency: a pairwise, not per-arm, condition]
\label{cor:pairwise}
Assume Theorem~\ref{thm:main-bound}(i) (or (ii) with $\rho_k \to 0$ and
$\bar v_k \to 0$), and fix the design shape as $C \to \infty$. Then:
\begin{enumerate}
  \item[(a)] \emph{Sufficient.} If $D_k^\dagger = \Delta_k + \mathrm{bias}(D_k) > 0$ for
    every $k \ne k^\star$, then $\mathbb{P}[\hat{k} = k^\star] \to 1$ (consistency in
    probability).
  \item[(b)] \emph{Failure.} If $D_k^\dagger < 0$ for some $k \ne k^\star$, then
    $\mathbb{P}[\hat{k} = k^\star] \to 0$ for that pair's comparison: selection converges
    to a wrong arm.
  \item[(c)] \emph{Boundary.} If $D_k^\dagger = 0$ for some $k$ (a limiting tie), neither
    conclusion follows from this argument: whether $\hat{k} = k^\star$ depends on the
    $O(\sqrt{v})$ fluctuations and on the tie-breaking rule, or, for a deterministic
    limiting tie, on how the tie is resolved. A tie resolved in favour of the true
    winner is consistent without a strictly positive gap, so strict positivity is
    \emph{not} necessary; it is only what this argument can guarantee.
\end{enumerate}
In particular, consistency does \emph{not} require $\sigma^2_{\mathrm{extrap}, k} = 0$
for every $k$: two arms whose individual biases are both nonzero but (nearly) equal at
$s^\star$ have $\mathrm{bias}(D_k) \approx 0$ even though neither
$\sigma^2_{\mathrm{extrap}, k}$ nor $\sigma^2_{\mathrm{extrap}, k^\star}$ is small. This
is the precise sense in which the \emph{pairwise} form is the operative one for the
decision problem, matching P1-06 step 4 / P1-07's empirical finding that the pairwise
bound is tighter than the marginal one (median tightness ratio $\approx 20.6$ vs.\
$\approx 24.1$ on real DataDecide cells) -- correlated errors between two arms partially
cancel in $D_k$ in exactly the cases this corollary identifies as consistency-preserving.
\end{corollary}

\begin{proof}
Convergence is in \emph{probability}, not almost surely: $v_k(C) \to 0$ alone does not
give almost-sure convergence, which would additionally need summable tail probabilities
along the sequence of designs. By the sub-Gaussian tail underlying
Theorem~\ref{thm:main-bound}, $\mathbb{P}[|D_k - D_k^\dagger| > \epsilon] \le
2\exp(-\epsilon^2 / (2(v_k + v_{k^\star})))$ (exactly in case~(i)), which tends to $0$
for every fixed $\epsilon > 0$ as $v \to 0$. Under~(a), taking $\epsilon = \min_k
D_k^\dagger / 2$ gives $\mathbb{P}[D_k > 0 \ \forall k] \to 1$, i.e.\ $\hat{k} = k^\star$
with probability tending to $1$. Under~(b), taking $\epsilon = |D_k^\dagger|/2$ gives
$\mathbb{P}[D_k < 0] \to 1$ for that pair. Case~(c) is not resolved by the argument.
\end{proof}

\subsection{Corollary 2 -- single-scale as a degenerate case}
\label{sec:corollary2}

\begin{corollary}[\texttt{ConstantExtrapolator} recovers the single-scale baseline]
\label{cor:single-scale}
With $g(\theta, s) = \theta$ (constant, $p=1$; \texttt{src/pdt/scaling/fitters.py}'s
\texttt{ConstantExtrapolator}) fit on a single proxy scale $S_{\mathrm{fit}} = \{s_p\}$:
$v_k(C) = \sigma^2(s_p) / n_{\mathrm{replicates}}$ (the ordinary sample-mean variance at
$s_p$), and
\[
  \sigma^2_{\mathrm{extrap}, k} = \big( \mu_k(s_p) - \mu_k(s^\star) \big)^2,
\]
the squared \emph{ordering distortion} of substituting the proxy scale's value for the
target's -- not an extrapolation bias in the usual sense (there is no fitted trend),
but the formal statement of what "single-scale accepts a fixed proxy gap instead of
paying an extrapolation bias" means, and why P1-03/04's single-scale-frontier
comparison is a special case of this same theorem, not a different one.
\end{corollary}

\begin{proof}
With $p=1$ and one fitting scale, $\hat\theta = \bar{y}(s_p)$ (the sample mean at
$s_p$) exactly; $J \equiv 1$ (constant, since $g$ doesn't depend on $s$), so
$v_k(C) = \Sigma_\theta(C) = \sigma^2(s_p)/n_{\mathrm{replicates}}$ directly. The
population projection of a single point is that point itself:
$\theta_k^\dagger(\{s_p\}) = \mu_k(s_p)$ exactly (there is no other scale to average
over), so $g(\theta_k^\dagger, s^\star) = \theta_k^\dagger = \mu_k(s_p)$, giving
$\sigma^2_{\mathrm{extrap}, k} = (\mu_k(s_p) - \mu_k(s^\star))^2$ directly from
Definition~\ref{def:sigma2-extrap}. Mechanical substitution into the general
definitions; no new argument.
\end{proof}

\subsection{Numerical certificate}
\label{sec:theorem1-certificate}

\texttt{tests/theory/test\_theorem1.py} simulates 5000 random instances (3--6 arms,
random linear-in-$\theta$ ground truth, random bounded sinusoidal misspecification
$h_k$ with $\sup_s |h_k(s)| = \eta_k$ realized exactly, random noise level not tied
monotonically to scale per Remark~\ref{rem:noise-not-monotone}) against
Theorem~\ref{thm:main-bound}'s bound (the corrected, gap-reduction form). Every
population quantity ($\theta_k^\dagger$, $\sigma^2_{\mathrm{extrap},k}$, $v_k$) is
computed in exact closed form (ordinary least squares on a linear-in-$\theta$ family),
so this certificate isolates Steps~\ref{sec:theorem1-step2}--\ref{sec:theorem1-step4}
of the proof from Step~\ref{sec:theorem1-step1}'s nonlinear-$g$ caveat. It is a check of
Theorem~\ref{thm:main-bound}(i) only, in a linear Gaussian setting; it is \emph{not}
evidence for the conditional nonlinear statement~(ii), whose hypothesis (H) it never
exercises.

\textbf{Result: 0 violations across all 5000 instances.} Of these, 1801 instances had a
theoretical bound large enough ($\ge 10^{-3}$) for 20000 Monte-Carlo trials to say
anything statistically meaningful about (checked via a one-sided Clopper-Pearson upper
confidence bound at $\alpha = 10^{-5}$ on the true error rate); the remaining 3199 had
bounds too small for Monte Carlo to resolve at this budget and are instead covered by
the closed-form argument in Steps~\ref{sec:theorem1-step2}--\ref{sec:theorem1-step4}
directly. Among the 1320 MC-checkable instances with at least one observed error,
tightness ratio (bound / empirical rate) ranged from $1.0$ (minimum -- the bound was
achieved essentially exactly, i.e.\ not looser than necessary, in the instance closest
to the worst case) to $2.0 \times 10^4$ (maximum, a design where realized bias signs
happened to favor correct selection far more than the worst-case treatment credits),
median $7.1$, IQR $[2.6, 41]$.
